\section*{Chapitre \ref{ch:m2:interest-rate-swaps} --- Swaps de taux}

\begin{solution}{exo:m2:interest-rate-swaps:1}
$100\,000\,000 \times 0.0405 \times 365/360 = \text{USD}~4\,106\,250$. Il reçoit le SOFR
composé sur les mêmes 365 jours, actual/360, sur USD~100 millions, versé deux
jours ouvrés après la fin de la période.
\end{solution}

\begin{solution}{exo:m2:interest-rate-swaps:2}
Recevoir le fixe revient à être long en duration, comme détenir une obligation à trente ans financée au
taux au jour le jour~: il perd quand les taux montent, car les engagements
(qu'il couvre) perdent eux aussi de la valeur.
\end{solution}

\begin{solution}{exo:m2:interest-rate-swaps:3}
$4.05 - 4.20 = -15$ points de base. Avant 2008, lorsque les \omterm{def:m2:interest-rate-swaps:swap}{jambes variables} versaient un
taux interbancaire porteur du risque de crédit bancaire, un spread négatif aurait signifié que
le marché valorisait l'emprunt à terme des banques en dessous de celui de l'État, ce qui ne
se produisait pas~; les \omterm{def:m2:interest-rate-swaps:spread}{spreads de swap} étaient positifs.
\end{solution}

\begin{solution}{exo:m2:interest-rate-swaps:4}
$N A \times 10^{-4}$ est la sensibilité à une variation du taux fixe propre au swap
à facteurs d'actualisation de la courbe inchangés~; choquer l'entrée au pair et reconstruire
la courbe modifie aussi les facteurs d'actualisation qui valorisent la \omterm{def:m2:interest-rate-swaps:swap}{jambe fixe}, de sorte que les deux
valeurs diffèrent d'un terme du second ordre (USD~16 sur 82\,000). Couvrir avec le nombre reconstruit et segmenté~:
c'est ce que font réellement les prix des instruments de couverture.
\end{solution}

\begin{solution}{exo:m2:interest-rate-swaps:5}
Un \omterm{def:m2:interest-rate-swaps:payer}{swap payeur} 5y5y gagne USD~82\,651 par point de base sur l'entrée à dix ans et perd
USD~45\,394 sur l'entrée à cinq ans~: il est long du taux à dix ans et court sur celui à cinq ans,
puisqu'il équivaut au swap à dix ans moins le swap à cinq ans. Couverture~: recevoir le fixe sur
un swap au pair à dix ans ayant le même DV01 à dix ans (USD~100.5 millions, puisqu'un
swap au pair à dix ans de USD~100 millions affiche 82\,262 sur ce pilier) et payer le fixe
sur un swap au pair à cinq ans affichant 45\,394 (USD~100.3 millions, soit 45\,278 pour 100
millions), ne laissant que de faibles résidus sur les autres piliers.
\end{solution}

\begin{solution}{exo:m2:interest-rate-swaps:6}
Notionnel brut de USD~300 millions~; chaque dealer paie le fixe sur un swap et reçoit
sur un autre aux conditions identiques, de sorte que son risque net est nul. La \omterm{def:m2:interest-rate-swaps:sef}{compression de portefeuille} résilie
les trois swaps~: le notionnel brut tombe à zéro, de même que les marges et le capital
que les trois transactions consommaient, sans modifier le risque de quiconque.
\end{solution}

\begin{solution}{exo:m2:interest-rate-swaps:7}
\omterm{def:m2:interest-rate-swaps:par}{Taux de swap au pair} à quatre ans 3.8425~\%, à huit ans 3.9743~\%, facteur d'actualisation à dix ans
0.666773. Les \omterm{def:m2:interest-rate-swaps:par}{taux de swap au pair} intermédiaires dépendent de l'interpolation~; les données
d'entrée, elles, n'en dépendent pas.
\end{solution}

\begin{solution}{exo:m2:interest-rate-swaps:8}
Un \omterm{def:m2:interest-rate-swaps:ois}{swap de taux au jour le jour} ne comporte presque aucun risque de crédit sur sa \omterm{def:m2:interest-rate-swaps:swap}{jambe variable}, et son
risque de contrepartie est éliminé par la compensation~: son taux n'est pas le coût d'emprunt
d'une banque. Un spread négatif signifie que recevoir le fixe sur un swap revient moins cher à porter
que la détention d'un Treasury, lequel doit être financé sur un bilan coûteux, et que la demande
de long terme pour recevoir le fixe (fonds de pension) excède ce que les dealers
sont prêts à absorber. Il s'agit d'un prix du bilan et de la demande, non du crédit.
\end{solution}

\begin{solution}{pb:m2:interest-rate-swaps:1}
\textbf{1.} 4.05~\%~; taux tout compris $4.05 + 1.50 = 5.55\%$.
\textbf{2.} $+\text{USD}~164\,474$ par point de base~: le payeur gagne quand les taux montent.
\textbf{3.} Nulle au \omterm{def:m2:interest-rate-swaps:par}{taux de swap au pair} de milieu de marché.
\textbf{4.} Le seul pilier à dix ans.
\textbf{5.} Les swaps en dollars démarrent deux jours ouvrés après la transaction~: le 25 septembre
2026 était un vendredi.
\textbf{6.} $+\text{USD}~8.02$ millions.
\textbf{7.} $-\text{USD}~8.44$ millions.
\textbf{8.} Un \omterm{def:m2:interest-rate-swaps:payer}{swap payeur} est court sur une obligation à taux fixe et long sur une obligation à taux variable~; il
est court en convexité, si bien qu'il gagne moins à la hausse qu'il ne perd à une baisse égale.
\textbf{9.} Chocs de $-20.0$, $-16.1$, $-12.2$, $-4.4$, $+3.3$ et $+15.0$ points de base à
un, deux, trois, cinq, sept et dix ans~; valeur $+\text{USD}~2.46$ millions.
\textbf{10.} Son risque repose entièrement sur le pilier à dix ans, qui a bougé de 15 points
de base~: $15 \times 164\,474 = 2.47$ millions, moins un petit terme de convexité.
\textbf{11.} Non~: la perte du swap est compensée par la moindre valeur de la dette à taux
fixe qu'il a synthétiquement créée~; son coût d'intérêt reste de 5.55~\%. La perte
correspond à celle d'avoir fixé à 4.05~\% alors que le marché fixerait désormais plus bas.
\textbf{12.} Comptabiliser les variations de valeur du swap en regard de celles de l'élément couvert, de sorte
que le résultat publié ne fluctue pas avec les taux lorsque la position économique est
couverte~; les trésoriers y tiennent, car des résultats volatils coûtent cher à expliquer.
\textbf{13.} Deux points de base pendant dix ans sur son DV01~: environ $2 \times 164\,474 =
\text{USD}~328\,947$ de valeur actualisée.
\textbf{14.} Chacune est exposée au défaut de l'autre lorsque le swap a de la valeur pour
elle~; il est géré par le collatéral prévu dans un credit support annex (CSA), par la compensation centrale
pour les dealers, et pour les clients entreprises par des lignes de crédit et des frais
intégrés au prix.
\textbf{15.} Les intérêts des années six à dix, qui redeviendraient variables.
\textbf{16.} Il conserve la souplesse et la tarification du prêt bancaire, peut être ajusté ou
dénoué séparément, et fixe le taux sans la documentation d'une émission obligataire
et la base d'investisseurs.
\textbf{17.} En recevant le fixe sur le marché interdealer ou sur un SEF, ou avec
des contrats à terme sur Treasury et des obligations, au même DV01 et sur le même pilier, puis en gérant le
risque résiduel de courbe dans son portefeuille.
\textbf{18.} L'entreprise a une dette à taux variable et veut fixer ce qu'elle paie~: elle paie
le fixe. Un fonds de pension a des engagements longs à taux fixe et veut des actifs qui gagnent
quand les taux baissent~: il reçoit le fixe.
\textbf{19.} Résultat retenu~: \emph{la couverture d'entreprise} fixe à 4.05~\% (5.55~\% tout compris),
avec un DV01 de USD~164\,474 par point de base, et vaut USD~2.46 millions après la
pentification.
\textbf{20.} Un taux fixe contre un taux variable, sur un notionnel qui ne change jamais
de mains.
\end{solution}

\begin{solution}{iq:m2:interest-rate-swaps:1}
Un échange d'intérêts à taux fixe contre des intérêts à taux variable sur un montant
notionnel, selon un échéancier. Une entreprise dont la dette est à taux variable paie le fixe pour
verrouiller son coût~; une entreprise dont la dette est à taux fixe peut recevoir le fixe pour profiter de la baisse des taux~;
un investisseur s'en sert pour ajouter ou retirer de la duration sans acheter ni vendre d'obligations.
\iqlookfor{l'échange de flux de trésorerie et un usage concret.}
\end{solution}

\begin{solution}{iq:m2:interest-rate-swaps:2}
L'intérêt au jour le jour composé sur une période se reproduit en renouvelant un dépôt,
donc dans un monde à courbe unique, le paiement de chaque période vaut $P(t_{i-1}) -
P(t_i)$, et la jambe entière $P(t_0) - P(t_n)$. Avec des courbes d'actualisation et
de projection distinctes, projeter le taux composé à partir de la courbe forward et
actualiser sur la courbe de collatéral.
\iqlookfor{l'argument télescopique et la prise en compte de la tarification multi-courbes.}
\end{solution}

\begin{solution}{iq:m2:interest-rate-swaps:3}
Choisir les instruments (dépôts ou contrats à terme pour le court terme, swaps au-delà), des piliers
placés à leurs échéances, et une interpolation (des logarithmes des facteurs d'actualisation, des taux
zéro ou des taux forward)~; résoudre pour chaque pilier à tour de rôle pour que son instrument soit
retarifé, avec une méthode de recherche de racine, ou toutes d'un coup. Test~: chaque donnée d'entrée est
retarifée à la tolérance près~; les taux forward ne présentent pas d'oscillation parasite~; modifier une
donnée d'entrée ne modifie la courbe que localement~; la courbe est stable à l'ordre des données~; les
courbes publiées sont reproduites.
\iqlookfor{les choix (instruments, interpolation) et une liste de tests.}
\end{solution}

\begin{solution}{iq:m2:interest-rate-swaps:4}
Détenir un Treasury consomme du bilan et doit être financé en repo~; recevoir le fixe sur un swap
ne demande que des marges. La demande de duration de long terme, notamment de la part des fonds de
pension, se satisfait à meilleur coût avec les swaps, et les bilans des dealers sont trop coûteux
pour arbitrer l'écart. La \omterm{def:m2:interest-rate-swaps:swap}{jambe variable}, hier indexée sur le LIBOR et aujourd'hui sur le SOFR
ou un taux au jour le jour, porte peu de risque de crédit, si bien que rien ne pousse les taux de swap
au-dessus des rendements des emprunts d'État.
\iqlookfor{bilan et demande, non le crédit.}
\end{solution}

\begin{solution}{iq:m2:interest-rate-swaps:5}
Recevoir un 5y5y revient à être acheteur de la baisse du taux à dix ans et vendeur de celle du taux à cinq ans~:
il s'agit du \omterm{def:m2:interest-rate-swaps:payer}{swap receveur} à dix ans moins le \omterm{def:m2:interest-rate-swaps:payer}{swap receveur} à cinq ans. Se couvrir en payant le fixe
sur un swap à dix ans au pair et en recevant sur un swap à cinq ans au pair, dans les DV01
correspondantes~; le résidu est une exposition à la forme de la courbe entre les
piliers, qui dépend de l'interpolation.
\iqlookfor{la décomposition en deux swaps au pair.}
\end{solution}

\begin{solution}{iq:m2:interest-rate-swaps:6}
Non-linéarité~: le choc parallèle n'est pas la somme des chocs locaux pour une
position convexe (termes du second ordre)~; une interpolation qui couple les piliers~;
des chocs sur les \omterm{def:m2:government-bonds:zero}{taux au pair} dans un calcul et sur les taux zéro dans un autre~; des réglages
de reconstruction de courbe ou des tolérances qui varient d'un calcul à l'autre~; des instruments hors
de l'ensemble des piliers (extrapolation)~; une courbe d'actualisation choquée dans un calcul et non dans
l'autre~; le bruit numérique provenant d'une tolérance trop lâche du chercheur de racine.
\iqlookfor{effets du second ordre et incohérences d'implémentation.}
\end{solution}
